\chapter{Literature review: markets and agency}\label{ch:litrev-agency}
% Assembled from the related-work sections of the papers in Part I
% (litrev/market-rl.tex, litrev/agency.tex). Lifted text is verbatim;
% composed glue is draft-marked.

\begin{draftblock}
\noindent The two chapters of this part draw on four bodies of work, which we review in turn: reinforcement learners constructed as internal markets, market models of bounded rationality, analogies between neural networks and markets, and the economics of aggregating preferences as applied recently to alignment problems.
\end{draftblock}

\section{Market-based reinforcement learning}\label{litrev:sec:marketrl}
% \label{mba/sec:related} kept from the source paper
\label{mba/sec:related}
\textbf{Market-based RL.} The majority of early work in this area has focused on market-based \emph{reinforcement learning} (RL) algorithms. The pioneering work in this domain consists of Holland's \emph{Learning Classifier Systems} or ``bucket brigade'' \citep{hollandPropertiesBucketBrigade1985a} in rule-based systems, where condition-action agents (``classifiers'') bid to post messages (which can be actions described in some language) onto a global message board. Improvements to this paradigm were made by Schmidhuber \citep{schmidhuberEvolutionaryPrinciplesSelfreferential1987, schmidhuberLocalLearningAlgorithm1989} who allowed agents to determine their own bids and imposed credit conservation, and by Baum \citep{baumModelIntelligenceEconomy1999} who further strictly enforced property rights, resulting in a much more familiar set-up called the ``Hayek Machine''. The Hayek machine, which can be applied to any Markov Decision Process (MDP), was subsequently extended to POMDPs by \citet{kweeMarketBasedReinforcementLearning2001} by adding external memory, and more recently \citep{changDecentralizedReinforcementLearning2020} modified the framework to use Vickrey auctions and prove that a Nash equilibrium of the market produces a globally optimal policy.\todo{is this much description necessary?}

\draft{All of these algorithms auction a single ``right to act'' at each step. \Cref{ch:market-rl} instead splits the state into separately traded goods, and \cref{ch:agency} develops the corresponding theory of prices.}

\section{Markets and bounded rationality}\label{litrev:sec:bounded}
\textbf{Bounded rationality and markets.} Other, more recent work in this area includes: \emph{Logical Induction} \citep{garrabrantLogicalInduction2020}, an algorithm that assigns probabilities to mathematical sentences based on their prices in a prediction market that (roughly speaking) pays off when a sentence is proven; and the \emph{Boundedly Rational Inductive Agent} (BRIA) \citep{oesterheldTheoryBoundedInductive2023} which solves finite decision problems by assigning it to the highest-bidding trader, similar to in market-based RL. The key insight of these works is that markets are useful for modeling \emph{boundedly rational agents}. The main results of each work -- the fact that the logical inductor cannot by dominated by any polynomial-time trader, and the ``boundedly rational inductive agent criterion'' in the latter paper -- are specific and precise formulations of the Efficient Market Hypothesis \citep{neymanAlgorithmicBayesianEpistemology2024}.\todo{remove last sentence?}

\section{Markets, neural networks and learning dynamics}\label{litrev:sec:nn}
\textbf{Markets and neural networks.} A specific equivalence between classifier systems and neural networks has been studied in the \emph{Neural bucket brigade} \citep{davisMappingClassifierSystems1988, schmidhuberLocalLearningAlgorithm1989}, although this does not consider backpropagation. A suggestive analogy between markets and backpropagation is discussed in \citet{wentworthCompetitiveMarketsDistributed2018}, though only for the case of a strictly sequential, unit-width market like that in \cref{mba:def:deep}. \draft{\Cref{ch:market-rl} makes this more precise and generalizes it to the wide markets of \cref{mba:def:goods}, and \cref{ch:agency} turns it into an exact equivalence.}\todo{maybe mention Agoric Computation, BitTensor and market-based scaffolding for LLMs? see comments}

\draft{From the learning dynamics perspective, the replicator equation of evolutionary game theory is known to be a form of Bayesian inference, where fitness can be interpreted as log-likelihood \citep{harperReplicatorEquationInference2009,baezInformationGeometryPart92012}; \cref{ch:agency} exploits this correspondence to interpret rewards and revenues as evidence. Proportional-response dynamics, in which each seller allocates its goods in proportion to the payments it receives, are known to converge to equilibrium in Fisher markets and can be derived as mirror descent \citep{birnbaum2011}. Analogies between economic equilibrium and statistical mechanics, which motivate some of the language of \cref{ch:agency}, are surveyed by \citet{liuEquilibriumThermoeconomics}.}

\section{Markets of AI agents and economics for alignment}\label{litrev:sec:econ-alignment}
\begin{draftblock}
Other literature takes a different approach and studies AI systems as participants in markets rather than as markets themselves. Early proposals for ``agoric'' computation envisaged software objects trading with each other for computational resources \citep{millerMarketsComputationAgoric1988}. More recent examples include decentralized networks that reward machine-learning models for their contributions \citep{raoBitTensorPeertoPeerIntelligence2021}, realistic simulations of economies in which language models are the agents \citep{liEconAgentLargeLanguage2024a}, and the use of simulated economies to inform tax policy \citep{zhengAIEconomistOptimal2021}. Most related to this thesis are recent works that take an economic perspective on alignment itself: \citet{conitzerPositionSocialChoice2024} and \citet{geAxiomsAIAlignment2024a} apply social choice theory to alignment, in particular to the aggregation of human feedback, and \citet{bergemannMechanismDesignAlignment2026} frame alignment as a problem of mechanism design. These works take the reward as given and ask how to elicit or aggregate it, whereas in \cref{ch:agency} the question is what the reward is, and the answer is given in terms of the prices in an economy.
\end{draftblock}
